\subsection{Heuristic-Guided Deep Reinforcement Learning (HuDRL)}

To address the high-dimensional resource allocation problem inherent in municipal budgeting, this study develops a \textit{Heuristic-Guided Deep Reinforcement Learning (HuDRL)} framework using the PyTorch-based \texttt{Stable-Baselines3} library \citep{Raffin2021}. While traditional optimization models often struggle with the dynamic, non-linear human behaviors found in waste management \citep{Saif2022}, Deep Proximal Policy Optimization (PPO) provides a robust architecture to approximate an optimal, adaptive governance policy $\pi_\theta(a|s)$ within a continuous action space. By coupling the DRL agent with the Agent-Based Model (ABM), the system evaluates the intricate socio-environmental feedback loops that dictate policy success or failure \citep{Jimenez2025, TianReview2024}.

\subsubsection{Proximal Policy Optimization Architecture}

To operationalize this policy engine, the simulation employs a custom Actor-Critic neural network architecture. PPO is specifically utilized for its mathematical stability; it prevents the agent from making erratic, catastrophic shifts in policy during training by strictly clipping the policy gradient update \citep{Raffin2021}. 

To effectively process the multi-modal state vector, the artificial intelligence is split into two distinct but cooperative Multi-Layer Perceptrons (MLPs):
\begin{itemize}
    \item \textbf{The Actor Network (The Decision Maker):} This network functions as the executive authority. It evaluates the current socio-environmental state of the city and outputs a specific action—deciding exactly how to divide the quarterly budget across different barangays and intervention types.
    \item \textbf{The Critic Network (The Evaluator):} Functioning as the internal advisor, this network evaluates the current state and predicts the long-term expected success (the Value) of the Actor's proposed strategy. By continuously grading the Actor's choices, the Critic guides the learning process toward sustainable governance rather than short-term fixes \citep{Mousavi2021}.
\end{itemize}

Both the Actor and Critic networks utilize two fully connected hidden layers comprised of 64 neurons each. This depth achieves a necessary computational balance: lightweight enough for rapid simulation, yet dense enough to process complex variables. These layers employ Rectified Linear Unit (ReLU) activation functions, which are highly effective at capturing the non-linear "tipping points" of human behavior—allowing the network to mathematically understand that minor increases in funding may yield zero behavioral change until a critical threshold is violently crossed \citep{Centola2018}.

\subsubsection{Markov Decision Process (MDP) Formulation}

\paragraph{State Representation ($S_t$)}
The state space $S_t \in \mathbb{R}^{17}$ constitutes the comprehensive observation vector accessible to the LGU Mayor Agent at the onset of each quarterly decision epoch $t$, formalized as:
\begin{equation}
    S_t = [\mathbf{C}_{1...7}, \mathbf{A}_{1...7}, B_{\mathrm{norm}}, T_{\mathrm{norm}}, P_{\mathrm{cap}}]
    \label{eq:state_vector}
\end{equation}
\addequation{17-Dimensional State Vector Representation}

\noindent This vector provides a holistic, normalized snapshot of the municipality. It includes the real-time compliance rate ($\mathbf{C}$) and average intrinsic attitude ($\mathbf{A}$) for all seven constituent barangays. It tracks the normalized remaining balance of the annual fund ($B_{\mathrm{norm}}$) and the current temporal progress of the simulation ($T_{\mathrm{norm}}$). Crucially, it incorporates Political Capital ($P_{\mathrm{cap}}$), which acts as a dynamic constraint tracking public approval.

\paragraph{Action Space ($A_t$) and Budget Constraint}
The continuous action space ($A_t \in \mathbb{R}^{21}$) dictates the simultaneous allocation of resources across three levers (IEC, Enforcement, Incentives) for all seven barangays. In municipal governance, the most absolute constraint is the finite nature of public funds. To mathematically enforce this without clipping the neural network's gradients, the raw output vector ($A_{\mathrm{raw}}$) is processed through a \textit{Temperature-Scaled Softmax} function ($\tau = 10.0$):
\begin{equation}
   A_{\mathrm{final}, i} = \frac{\exp(\tau \cdot A_{\mathrm{raw}, i})}{\sum_{j=1}^{21} \exp(\tau \cdot A_{\mathrm{raw}, j})} \cdot B_{\mathrm{Quarterly}}
   \label{eq:softmax_constraint}
\end{equation}
\addequation{Temperature-Scaled Softmax Budget Constraint}
This exponential scaling aggressively amplifies the AI's highest confidences, forcing it to make polarized, decisive investments rather than spreading funds too thinly. The normalization mathematically guarantees that the sum of all allocated funds exactly equals the Quarterly Budget Cap ($\sum A_{\mathrm{final}} \le B_{\mathrm{Quarterly}}$) \citep{Abdallah2021}.

\paragraph{Multi-Objective Reward Function ($R_{\mathrm{total}}$)}
The agent is guided by a composite reward function designed to balance environmental outcomes, fiscal responsibility, and social acceptability \citep{Hertweck2023Values}:
\begin{equation}
    R_{\mathrm{total}} = R_{\mathrm{Base}} + R_{\mathrm{Jackpot}} + R_{\mathrm{Focus}} - P_{\mathrm{Waste}} \pm P_{\mathrm{Cap}}
    \label{eq:total_reward}
\end{equation}
\addequation{Composite Multi-Objective Reward Function}
The AI receives a base reward for global compliance ($R_{\mathrm{Base}}$). However, to solve the \textit{Sparse Reward Problem}, heuristic "breadcrumbs" are added. The agent receives a massive $R_{\mathrm{Jackpot}}$ (+200) if global compliance breaches the 70\% threshold. It earns a Focus bonus ($R_{\mathrm{Focus}}$) for heavily funding the worst-performing barangay, and suffers an Efficiency Penalty ($P_{\mathrm{Waste}}$) if it allocates excess funds to a barangay that has already crossed the 70\% target.

\subsubsection{Addressing Behavioral Inertia: The Triage Heuristic}
In a standard, unguided PPO environment, an agent explores this massive action space at random. Because overcoming community-wide behavioral inertia requires a critical mass of intervention to trigger a ``Social Tipping Point'' \citep{Nyborg2024}, random exploration rarely yields a sudden increase in compliance, causing agents to converge to a ``Status Quo'' local optimum.

To overcome this, the HuDRL framework integrates a domain-specific "Triage Logic" heuristic before the Softmax normalization step. This acts as a saliency filter \citep{Alphonso2024}, forcibly categorizing barangays into strategic states.

\begin{algorithm}[H]
\caption{The Mayor Agent's Heuristic Triage Logic}
\label{alg:hudrl_triage_english}
\begin{algorithmic}[1]
\State \textbf{Input Data:}
\State \parbox[t]{0.9\linewidth}{$\rightarrow$ The AI's raw proposed action vector (21 values).\strut}
\State \parbox[t]{0.9\linewidth}{$\rightarrow$ Real-time compliance data for all 7 Barangays.\strut}
\Statex
\State \textbf{Step 1: The 70\% Graduation Rule (Identify Solved Areas)}
\For{Each Barangay in the Municipality}
    \If{Barangay Compliance Rate $\ge$ 70\%}
        \State \parbox[t]{0.85\linewidth}{$\rightarrow$ \textit{Budget Slash:} Multiply the AI's proposed budget allocation for this specific barangay by $0.01$ (reducing it by 99\%).\strut}
        \State \parbox[t]{0.85\linewidth}{$\rightarrow$ \textit{Rationale:} This prevents the AI from becoming a "Victim of Success" and wasting municipal funds maintaining an area where social norms have already locked in the behavior.\strut}
    \EndIf
\EndFor
\Statex
\State \textbf{Step 2: Calculate Final Budget Distribution}
\State \parbox[t]{0.9\linewidth}{$\rightarrow$ Take the adjusted raw values and pass them through the Temperature-Scaled Softmax equation.\strut}
\State \parbox[t]{0.9\linewidth}{$\rightarrow$ This automatically redirects the funds stripped from the "Graduated" barangays and funnels them heavily into the remaining struggling barangays.\strut}
\Statex
\State \textbf{Step 3: Evaluate \& Update (Reward Shaping)}
\If{The AI successfully focused its budget on a struggling area and increased its compliance}
    \State \parbox[t]{0.85\linewidth}{$\rightarrow$ Grant a massive scaling reward multiplier.\strut}
\EndIf
\If{The AI attempted to spend more than 10\% of the budget on a "Graduated" area}
    \State \parbox[t]{0.85\linewidth}{$\rightarrow$ Apply a negative penalty to teach the AI to abandon solved zones.\strut}
\EndIf
\end{algorithmic}
\end{algorithm}

This heuristic shaping forces the agent to explore the ``Sequential Saturation'' strategy—overwhelming a single problematic area with resources until social norms shift, before redirecting funds to the next \citep{Cheng2021}.